\documentclass[a4paper, 11 pt, fleqn]{article}

% Layout und Allgemeines
\usepackage[left = 2.5 cm, right = 2.5 cm, top = 2.5 cm, bottom = 2.5 cm ]{geometry} % Seitenränder
\usepackage[utf8]{inputenc} % Encoding
\usepackage{microtype} % verbesserter Randausgleich
\usepackage[onehalfspacing]{setspace} % Zeilenabstand 1,5
% Bibliographie
\usepackage[authoryear]{natbib}
% Sidenotes
\usepackage{sidenotes}
\usepackage{tikz}
\usetikzlibrary{calc}
\usepackage{makecell}
\usepackage{amsmath}
\usepackage{amssymb}
\usepackage{dsfont}
\usepackage{mathrsfs} % rsfs fonts in math mode
\usepackage{float}
\usepackage{multirow}
% Grafiken/Abbildungen
\usepackage{graphicx} 
\usepackage{subcaption}
\usepackage{pdflscape}
\usepackage{hyperref}

% Bunter Text
\usepackage{xcolor}
\graphicspath{{images/}}
\DeclareGraphicsExtensions{.pdf,.png,.jpg}
\setcitestyle{aysep={}, yysep={;}}
\bibliographystyle{apalike}

\DeclareMathOperator*{\argmax}{arg\,max} % argmax
\DeclareMathOperator*{\argmin}{arg\,min} % argmin

\tikzset{
  solid node/.style={circle,draw,inner sep=1.2,fill=black},
  hollow node/.style={circle,draw,inner sep=1.2},
}

\begin{document}
\begin{titlepage}
	\noindent
	Otto-Friedrich-Universitätt Bamberg\\
	Fakult\"at f\"ur Sozial- und Wirtschaftswissenschaften\\
	Lehrstuhl f\"ur Volkswirtschaftslehre, insb. Finanzwissenschaft\\
	Prof. Dr. Florian Herold
	
	\begin{center}
		\vspace{2 cm}
		\Large{
			\textbf{
		Consumer Responses to Behaviour-Based Price Discrimination
		}}\\ 
		\vspace{0.5 cm}
		\normalsize{Datum der Abgabe: 21.04.2023}\\
		\vspace{1 cm}
		\includegraphics[scale=0.175]{graphics/logo.png}\\
		\vspace{1 cm}	
		\textbf{Masterarbeit \\
		im Studiengang \\
		European Economic Studies (EES) \\
		an der Fakultät Sozial- und Wirtschaftswissenschaften \\
		der Otto-Friedrich-Universität Bamberg}\\
	\end{center}
	
	\vspace{5cm}
	\flushright{
		Niklas Sebastian Dörner\\
		An der Weberei 3, 96047 Bamberg\\
		Matrikelnr. 1859018\\
		niklas-sebastian.doerner@stud.uni-bamberg.de\\
	}			
\end{titlepage}

% Verzeichnisse
\pagenumbering{Roman}
\newpage
\addcontentsline{toc}{section}{Table of Contents}
\tableofcontents
\listoffigures
\newpage

% \listoftables

\newpage

\pagenumbering{arabic}

% Endlich kommt mal Inhalt
\begin{abstract}
	The abstract
\end{abstract}
%
\section{Introduction} \label{sec:intro}
%
\citet{Choe.2018} is chosen due to its significant contribution to the literature. The model they develop is dynamic and
was innovative in the sense that endogenously created asymmetric information leads to multiple and asymmetric equilibria with
and without endogenous product choice even when firms are symmetric ex ante. The asymmetry in equlibrium depends on the foward-looking
tendency of firms (positive discount factor for the future). \textcolor{red}{Siehe Content-Reiter Citavi}

\subsection{Price Differentiation} \label{ssec:price-diff}
Types of price discrimination
%
\subsection{Spatial Competition}
Mostly Hotelling (incl. impossibility result $\to$ relevance for Choe) and maybe also the circle (see Belleflamme)
%
\section{\citet{Choe.2018}: BBPD and Spatial Competition in a Dynamic Model} \label{sec:choe-2018-intro}
%
This section provides an introduction to the model presented in \citet{Choe.2018}. The present a model that combines
behaviour-based price discrimination with spatial competition. A brief overview of the model is followed by a thorough
exposition below. The model is dynamic and comprises two periods and two firms.
In the first period, firms choose their location as they do in the standard Hotelling model with quadratic transportation costs
but also track all consumers that decide to purchase from them in this period. This creates informational asymmetries as
firms can only track customers that purchased from \textit{them}. Once chosen, locations are fixed for the entirety of the
game. Note that the competiton for location is also a competition for information. In the baseline model, this thus concludes the
spatial competition between firms. Customer's location in the Hotelling linear city over [0,\,1] is the only source of heterogeneity
in the set of consumers.\footnote{This is sufficient as a source of heterogeneity as variation in willingness to pay based on different
transportation costs is the only property required by the model.} In the second period, firms can exploit their private information by
discriminating against consumers based on their decisions from the first period:~firms offer `personalised pricing' (p.~5672) to
\textit{each} consumer who purchased from that firm in the first period. Firms also set a uniform `poaching' price (p.~5672) in
order to attract consumers who previously purchased from their competitor. Firms simultaneously choose the poaching price first,
after which they set the individualised prices for their existing customers. This price structure hinges on the ability of firms
to track their customers in the first period and endogenously create private information. Consumers then observe three prices
prior to their purchase decision in the second period: their personalised price and the two poaching prices; they only care about
the lowest price absent any loyalty or fairness concerns. The authors assume that the inclusion of ``bevioral elements to consumer
choice is likely to change [the] results substantially.'' (p.~5672) The model is dynamic since firms have multiple periods to act,
multiple decisions to make and interact with one another. Accordingly, firms intertemporally maximise their total profits over both
periods. The authors thus endow firms (and consumers) with positive discount factors. The discounting of future profits by firms
ensures the existance of asymmetric equilibria; if firms have no intertemporal prefence for second-period profits, equilibria are
symmetric irrespecive of consumers' discount factors.
%
\subsection{Baseline Model} \label{sec:choe-baseline} % based on Section 2 of \citet{Choe.2018}
The firm comprises two firms that compete in a Hotelling linear city with quadratic costs for location and consumers in a two-stage game
with perods $\tau = 1,\,2$. Consumers are uniformly distributed on [0,\,1] where consumer $x$ is simply located at $x$.
Consumers incur quadratic transportation costs. Firms have full knowledge of this setup. Furthermore, firms are assumed to have constant
marginal cost of production; this is~--~as is common in this strain of the literature~--~normalised to 0. Over the course of the game,
firms gather information and engage in price discrimination against consumers. Consumers are price takers and simply act as a source of
demand without having any notable actions to chose from.
%
\subsubsection{Firms}
\textcolor{red}{Grafiken erstellen, um Spiel aus Sicht der Firma anschaulicher darzustellen?}
Firms are indexed by $i = A,\,B$ and act as duopolists. Firms compete for location for $\tau=1$. Since firms can only gather relevant
information in the first period, the location game is at the same time a contest for (private) information. The outcome of this game
determines the distribution of endogenously created private information among the dupolists and their respective ability to engage in
price discrimination. Firms set a unform price $P^1_i$ in the first period and track their customers. For each consumer, a firm knows
whether it was a customer in the first period. \citet[p. 5672]{Choe.2018} describe the information gain from tracking as follows:
let $\mathscr{A}$ denote the set of consumers that are a customer of firm $A$ in $\tau=1$. After $\tau=1$, firm $A$ knows
$\forall x \in [0,\,1]$ whether $x\in \mathscr{A}$ and also the location~$x$ of consumer~$x$ if that consumer was indeed a customer
of firm~$A$ in the first period. This similarly holds with customer set~$\mathscr{B}$ and customers~$y$ for firm~$B$.
The model structure assumed by \citet{Choe.2018} in conjunction with the tracking mechanism
implies that both firms know the location of the indifferent consumer \textcolor{red}{$\hat{x}$ benennen ja/nein?} (marginally above (below) the rightmost (leftmost) consumer in their
set of first-period customers, depending on whether the firm served the customers to left (right) of the indifferent consumer) and can, with
certainty, assign non-customers their first-period supplier since there are only two firms and the world is linear. If the model were to include
$n > 3$ competing firms, setting an efficient poaching price would likely be a much more involved process for any given firm. \\
% Period 2
The pricing strategies available to firm $i$ in period $\tau=2$ are determined by the informational asymmetries that result from the outcome
of the location game. \citet[p. 5672]{Choe.2018} let firms set two different types of prices in period $\tau=2$. Based on the private
information obatained in $\tau=1$, firms engage in price discrimination in two disctinct ways. They\footnote{For brevity I assume the perspective
of firm $A$ in the following description of price setting.} discriminate \textit{between} first-period customers (all $x\in\mathscr{A}$)
and consumers that chose the competitor (all $x\notin\mathscr{A}$, i.e. all $x\in\mathscr{B}$ since the authors assume a duopoly) but also
\textit{within} $\mathscr{A}$. Firms set a uniform poaching price, e.g. $P_A(\mathscr{B})$, aimed at all consumers who chose their competitor
in the first period and an individualised price, e.g. $P_A(x)$. The poaching price is designed to incentivise consumers to switch to another
firm in $\tau=2$ while the individualised pricing schedule is designed to appropriate the highest possible surplus based on the information
obtained in the first period.
However, each consumer chooses among their personsalised price and the two poaching prices. Note that the uniform poaching price can be viewed
as primarily a function of set membership, while the individualised price offered to a given customer $x\in\mathscr{A}$ is a function of their
specific location along the [0,\,1]-line \textcolor{red}{Mit oder ohne Bindestrich?}.
Firms set their poaching price first, followed by the individualised prices (p. 5672). \textcolor{orange}{ibid? Autor ja/nein? Nur Seitenzahl?}
Discriminating between the two major groups of consumers is \textcolor{red}{first/second/third; nachschauen was Choe et al. schreiben, aber auch mal neben Definitionen legen}
price discrimination. Price discrimination against the consumers in $\mathscr{A}$ (or $\mathscr{B}$, respectively) is first-degree price discrimination. \\
% Discount factor
\citet[p. 5672]{Choe.2018} endow firms with a positive discount factor $\delta_f \in [0,\,1]$. The discount factor is identical for both
firms but may differ from the discount factor of the consumer population, $\delta_c$. The model produces asymmetric equilibria as long as
$\delta_f > 0$ (p. 5673). Firms thus face an intertemporal decision problem to maximise their total discounted profit
$\Pi_i = \pi^1_i + \delta_f\pi^2_i$ and compete for location (information) based on their pricing regimes.
%

%
\subsubsection{Consumers}
The model assumes a single set of homogenous consumers. All consumers have the same valuation $v$ for the good sold by the firms and only 
differ in their location within Hotellings linear city.
%


\subsection{Related Literature (tbd)}
\begin{itemize}
	\item p. 5672: ``we abstract away from these issues mainly because our aim is to clearly understand how changes in informational assumptions lead to different equilibria
	in otherwise the same model as Fudenberg and Tirole (2000)''
	\item die Autoren nennen 3 Referenzmodelle, mit denen sie sich vergleichen (ua FT2000), ggf gibt es noch andere; die im Artikel genannten sind relativ alt
\end{itemize}
%
\section{\citet{Armstrong.2019}: Discriminating Against Captive Consumers}
%
\subsection{Captive Consumers~--~Consumer Myopia}
%
\section{\textcolor{red}{Zwischenfazit: What we have learned so far, where we stand now and where we are going next}}
Preliminary conclusion: Model Components oder so ähnlich
%
\section{Extensions} \label{sec:extensions}
%
\subsection{Intersecting \citet{Choe.2018} and \citet{Armstrong.2019}}
%
\subsection{Consumer Responses: Signalling and Search}
%
\subsubsection{Signalling}
%
\subsubsection{Search}
%
\newpage
\pagenumbering{Roman}
% \setcounter{page}{3}
\bibliography{bibliography}

\newpage
\section*{Selbständigkeitserklärung}
Ich erkläre hiermit gem. \text{\S}6 Abs. 6 APO SoWi, dass ich die vorstehende Masterarbeit selbstständig verfasst, keine anderen
als die angebenen Quellen und Hilfsmittel benutzt und die wörtlich oder inhaltlich übernommenen Stellen als solche kenntlich
gmacht habe. Des Weiteren erkläre ich, dass die digitale Fassung der gedruckten Asufertigung der Masterarbeit ausnahmslos
in Inhalt und Wortlaut entspricht und zur Kenntnis genommen wurde, dass diese digitale Fassung, einer durch Software
unterstützten, anonymisierten Prüfung auf Plagiate unterzogen werden kann. \\
\vspace*{1 cm}
\\
\noindent\begin{tabular}{ll}
	\makebox[2.5in]{\hrulefill} & \makebox[2.5in]{\hrulefill}\\
	Niklas D\"orner & Ort, Datum\\[8ex]% adds space between the two sets of signatures
\end{tabular}
%
\end{document}